In semantic segmentation, IoU is often computed independently for each class and then averaged across classes, yielding the \textbf{mean Intersection over Union (mIoU)}.\cite{tensorflow_mean_iou_2026,muller_guideline_2022} This metric is particularly informative in multi-class settings because it penalizes both over-segmentation and missed regions.

Other commonly adopted metrics include \textbf{precision} and \textbf{recall}, defined at pixel level as follows:\cite{sklearn_precision_recall_2026}

\begin{equation}
Precision = \frac{TP}{TP + FP}
\end{equation}

\begin{equation}
Recall = \frac{TP}{TP + FN}
\end{equation}

Precision measures how many of the pixels predicted as positive are actually correct, while recall measures how many of the truly positive pixels are successfully retrieved by the model.\cite{sklearn_precision_recall_2026} Their harmonic mean is given by the \textbf{F1-score}, or equivalently the \textbf{Dice coefficient}, which is widely used in segmentation tasks.\cite{muller_guideline_2022}

\begin{equation}
F1 = Dice = \frac{2TP}{2TP + FP + FN}
\end{equation}

In binary segmentation problems or anomaly-detection settings, the \textbf{Area Under the Receiver Operating Characteristic Curve (AUROC)} can also be used. AUROC evaluates the model's ability to discriminate between positive and negative pixels across different decision thresholds.\cite{sklearn_roc_auc_2026} However, in strongly imbalanced segmentation tasks, overlap-based metrics such as IoU or Dice are often considered more informative.\cite{muller_guideline_2022}

For instance segmentation, COCO-style evaluation commonly reports \textbf{Average Precision (AP)} over mask IoU thresholds, while panoptic segmentation is commonly evaluated with \textbf{Panoptic Quality (PQ)}.\cite{detectron2_coco_evaluation_2020,kirillov_panoptic_2019}
These evaluation metrics provide a quantitative measure of the agreement between predicted segmentation masks and ground-truth annotations, enabling objective comparison between different models and approaches.
